\chapter{Conclusion and Future Work}
\label{ch:conclusion}

% Chapter introduction
This concluding chapter synthesizes the comprehensive research journey of the thesis, ``A Two-Tier Cognitive Architecture for Automated Threat Detection and Response Using Agentic AI'' (SENTINEL). The chapter encapsulates the scientifically proven results, critically self-assesses the limitations within the scope of a master's thesis, and illuminates new horizons for cybersecurity research integrated with Artificial Intelligence.

\section{Synthesis of Achieved Research Results}
The synthesis of the achieved research results confirms the successful resolution of the core research questions. Specifically, the tiered architecture combining the Tier-1 Welford filter and the Tier-2 LangGraph Agent successfully mitigates the latency bottlenecks and alert fatigue that plague conventional SOC workflows. In terms of AI system security, the empirical evaluations demonstrate the viability of Delimited Data Encapsulation and HMAC log chaining in thwarting adversarial manipulation risks—including prompt injections, delimiter smuggling, and data tampering—that target internal enterprise cognitive systems. Furthermore, the system successfully overcomes Zero-day and APT challenges through the synergy between historical alert tracking via Threat Memory and the hybrid knowledge repository utilizing Dual-RAG. This provides the agent with profound contextual awareness, enabling it to expose highly sophisticated, multi-stage infiltration campaigns. Finally, the validation rigor of these findings is established through robust 5D evaluations and standard statistical tests, such as McNemar's and Mann-Whitney U tests, performed on empirical datasets.

\section{Scientific and Practical Contributions}
This thesis delivers both scientific and practical contributions to the field of cybersecurity. In terms of scientific value, this research provides a highly modular architectural framework for integrating mathematical statistics with Large Language Models in cybersecurity, substantially fortifying the theoretical repository concerning LLM cognitive security, guardrails, and context isolation. Practically, the finalized SENTINEL prototype serves as a deployable core engine or an independent auxiliary tool for Small and Medium Businesses (SMBs) that lack the capital to sustain large-scale Tier-1 SOC analyst teams. By operating entirely offline, the system strictly maintains data sovereignty and adheres to Zero Trust privacy guidelines.

\section{Existing Limitations}
Despite the system's demonstrated superiority, this research acknowledges several key limitations that define its operational boundary. Regarding input log formats, the implementation was evaluated primarily on Network Flow logs and HTTP payloads, whereas real-world enterprise environments require the ingestion of diverse unstructured formats—such as Windows Event Logs, Linux Syslogs, and CloudTrail records—demanding more complex semantic parsing layers. In terms of hardware constraints, although inference speed is accelerated using semantic caching, analyzing entirely novel alerts still requires substantial LLM processing cycles, mandating significant capital expenditures for dedicated GPU hardware if deployed at high-throughput network scales. Finally, regarding physical mitigation design, the research deliberately restricted the agent to generating recommended actions—such as proposing firewall rule updates—rather than executing automated physical blocks on network routers. While this constraint prevents the risk of accidental Self-Denial of Service (Self-DoS), it limits the system's capacity for fully autonomous, real-time blocking.

\section{Future Development Directions}
Several future development directions emerge from these limitations to guide subsequent research. First, the integration of federated learning could expand SENTINEL into a decentralized architecture, enabling multiple SOC endpoints to share learned threat parameters and model weights regarding malware mutations without transmitting privacy-violating raw network logs. Second, the monolithic LangGraph agent could be decomposed into a Multi-Agent System (MAS) comprising specialized sub-agents, such as malware, network, and threat intelligence analysts, which debate the evidence before issuing a final verdict to enhance classification accuracy. Third, the posture of the system could be evaluated by constructing an automated Red Team agent to continuously generate sophisticated adversarial payloads targeting SENTINEL, facilitating autonomous self-evaluation and defensive upgrades.

\vspace{1cm}
\begin{flushright}
    \textit{Closing a modest endeavor, yet opening profound aspirations in the pursuit of securing the digital realm through the power of Agentic AI.}
\end{flushright}
